\documentclass[11pt]{article}
\usepackage{amsmath, amssymb, amsthm}
\usepackage{hyperref}
\usepackage[margin=1in]{geometry}

\newtheorem{theorem}{Theorem}
\newtheorem{definition}[theorem]{Definition}
\newtheorem{example}[theorem]{Example}
\newtheorem{remark}[theorem]{Remark}

\title{A Guide to Elliptic Curve Invariants:\\Understanding the Four Tables}
\author{}
\date{}

\begin{document}
\maketitle

\begin{abstract}
We explain four tables of arithmetic and modular data for elliptic curves over $\mathbb{Q}$ with conductor $N \leq 54$, aimed at mathematicians with general background but limited number theory experience. We develop the theory from basics: elliptic curves as algebraic groups, the Mordell-Weil theorem, modular forms and the Modularity Theorem, Galois representations, and why representation-theoretic data matters for understanding L-functions and arithmetic. All technical terms are defined before use.
\end{abstract}

\tableofcontents
\clearpage

\section{Introduction: What Are We Looking At?}

These four tables catalog various invariants of elliptic curves over $\mathbb{Q}$---smooth projective curves of genus 1 with a specified rational point. While geometrically simple objects, elliptic curves connect to deep areas of modern number theory: they admit a group structure (making them abelian varieties), they parametrize modular forms via the Modularity Theorem, and they carry rich Galois representations encoding arithmetic information.

The tables split this data into natural categories:
\begin{enumerate}
\item \textbf{Table 1}: Basic invariants---weight, torsion, rank
\item \textbf{Table 2}: Mordell-Weil generators and Galois representations
\item \textbf{Table 3}: Galois exceptional primes and modular groups
\item \textbf{Table 4}: Modular curves and L-functions
\end{enumerate}

We also include data for Ramanujan's $\Delta$ function, a weight 12 modular form, which serves as a conceptual bridge between elliptic curves and more general modular forms.

\section{Background: Elliptic Curves and Their Groups}

\subsection{Definition and Basic Structure}

\begin{definition}
An \textbf{elliptic curve} over $\mathbb{Q}$ is a smooth projective curve $E$ of genus 1 with a specified rational point $O \in E(\mathbb{Q})$.
\end{definition}

Such curves can always be written in \textbf{Weierstrass form}:
$$E: y^2 = x^3 + Ax + B$$
where $A, B \in \mathbb{Q}$ and $\Delta = -16(4A^3 + 27B^2) \neq 0$ (smoothness condition). The specified point is typically $O = [0:1:0]$ at infinity.

The miracle: $E(\mathbb{Q})$ forms an \textbf{abelian group} under a geometric addition law. Given two points $P, Q \in E(\mathbb{Q})$, the line through $P$ and $Q$ intersects $E$ in a third point $R$, and we define $P + Q$ to be the reflection of $R$ across the $x$-axis. The point $O$ serves as the identity.

We write $[n]P$ for the $n$-fold sum:
$$[n]P = \underbrace{P + P + \cdots + P}_{n \text{ times}}$$

\subsection{The $j$-Invariant}

\begin{definition}
The \textbf{$j$-invariant} of an elliptic curve $E: y^2 = x^3 + Ax + B$ is:
$$j(E) = 1728 \cdot \frac{4A^3}{4A^3 + 27B^2}$$
\end{definition}

This is a rational function that completely determines the isomorphism class of $E$ over an algebraically closed field.

\textbf{Key fact}: Two elliptic curves over $\overline{\mathbb{Q}}$ are isomorphic if and only if they have the same $j$-invariant. Over $\mathbb{Q}$, the situation is more subtle---curves with the same $j$-invariant are related by twists.

\subsection{The Conductor}

The \textbf{conductor} $N$ is a positive integer measuring the arithmetic complexity of the curve. For each prime $p$, we reduce the Weierstrass equation modulo $p$:

\begin{definition}
An elliptic curve $E$ has:
\begin{itemize}
\item \textbf{Good reduction} at $p$ if the reduced curve $\tilde{E}/\mathbb{F}_p$ is smooth (nonsingular)
\item \textbf{Bad reduction} at $p$ if $\tilde{E}/\mathbb{F}_p$ is singular (has a node or cusp)
\end{itemize}
\end{definition}

Bad reduction comes in two types:
\begin{itemize}
\item \textbf{Multiplicative reduction}: $\tilde{E}$ has a node (ordinary double point)
\item \textbf{Additive reduction}: $\tilde{E}$ has a cusp
\end{itemize}

The conductor encodes this information:
$$N = \prod_{p \text{ bad}} p^{f_p}$$
where the exponent $f_p$ is defined by:
$$f_p = \begin{cases}
0 & \text{good reduction} \\
1 & \text{multiplicative reduction} \\
2 + \delta_p & \text{additive reduction}
\end{cases}$$
and $\delta_p \geq 0$ measures the ``wildness'' of the additive reduction (related to valuations of the minimal discriminant and $j$-invariant at $p$).

\textbf{Interpretation}: The conductor is both a geometric invariant (measuring singularities of reduction) and an arithmetic invariant (appearing in functional equations of L-functions). Smaller conductor indicates ``better'' arithmetic behavior.

\subsection{The Cremona Label}

Elliptic curves are organized into \textbf{isogeny classes}---equivalence classes under isogenies (surjective morphisms with finite kernel). The \textbf{Cremona label} has the form $NXm$:
\begin{itemize}
\item $N$ = conductor
\item $X$ = letter (a, b, c, ...) indicating isogeny class
\item $m$ = number ordering curves in the class
\end{itemize}

Example: \texttt{11a1}, \texttt{11a2}, \texttt{11a3} are three curves of conductor 11 in the first isogeny class.

\section{Table 1: Basic Arithmetic Invariants}

\subsection{Column A: Curve Label}

As described above. The tables include all curves with $N \leq 54$ from the Cremona database, plus the special entry ``$\Delta$'' for Ramanujan's cusp form.

\subsection{Column B: Weight}

This refers to the \textbf{weight of the associated modular form}. To understand this, we need the Modularity Theorem.

A \textbf{modular form} of weight $k$ for a congruence subgroup $\Gamma \subseteq \text{SL}_2(\mathbb{Z})$ is a holomorphic function $f: \mathbb{H} \to \mathbb{C}$ on the upper half-plane satisfying:
$$f\left(\frac{az+b}{cz+d}\right) = (cz+d)^k f(z)$$
for all $\begin{pmatrix} a & b \\ c & d \end{pmatrix} \in \Gamma$, plus growth conditions at cusps.

A \textbf{cusp form} additionally vanishes at all cusps. The space of weight $k$ cusp forms for $\Gamma_0(N)$ is denoted $S_k(\Gamma_0(N))$.

\subsubsection{Newforms, Oldforms, and Levels}

Before stating the Modularity Theorem, we must clarify crucial terminology about levels and newforms.

\begin{definition}
For a modular form $f$, we say $f$ \textbf{has level $N$} or \textbf{appears at level $N$} if $f \in S_k(\Gamma_0(N))$.
\end{definition}

A key observation: if $f \in S_k(\Gamma_0(N))$ and $M$ is any multiple of $N$ (i.e., $N | M$), then $f \in S_k(\Gamma_0(M))$ as well. This is because $\Gamma_0(M) \subseteq \Gamma_0(N)$, so the transformation property for $\Gamma_0(N)$ automatically implies the transformation property for the smaller group $\Gamma_0(M)$.

This means: \textbf{a single modular form appears at infinitely many levels}---all multiples of some minimal level.

\begin{definition}
The space $S_k(\Gamma_0(N))$ decomposes as:
$$S_k(\Gamma_0(N)) = S_k^{\text{new}}(\Gamma_0(N)) \oplus S_k^{\text{old}}(\Gamma_0(N))$$
where:
\begin{itemize}
\item $S_k^{\text{new}}(\Gamma_0(N))$ is the space of \textbf{newforms} at level $N$---forms that do \emph{not} come from any proper divisor of $N$
\item $S_k^{\text{old}}(\Gamma_0(N))$ is the space of \textbf{oldforms} at level $N$---forms that already appeared at some level $M$ with $M | N$ and $M < N$
\end{itemize}
\end{definition}

More precisely:
$$S_k^{\text{old}}(\Gamma_0(N)) = \sum_{\substack{d | N \\ d < N}} S_k^{\text{new}}(\Gamma_0(d))$$

A form is ``new'' at level $N$ if it genuinely first appears at that level, not having come from a lower level.

\begin{definition}
For a newform $f \in S_k^{\text{new}}(\Gamma_0(N))$, we call $N$ the \textbf{conductor} of $f$ (as a modular form). This is the minimal level at which $f$ appears.
\end{definition}

\textbf{Key distinction}:
\begin{itemize}
\item A newform $f$ \textbf{has conductor $N$} (unique, minimal, intrinsic)
\item The same form $f$ \textbf{appears at levels $N, 2N, 3N, \ldots$} (infinitely many, as oldform for $M > N$)
\end{itemize}

\begin{theorem}[Minimality of Newform Level]
Let $f \in S_k^{\text{new}}(\Gamma_0(N))$ be a newform. Then:
\begin{enumerate}
\item $f \in S_k(\Gamma_0(M))$ if and only if $N | M$
\item $N$ is the unique minimal positive integer such that $f \in S_k(\Gamma_0(N))$
\end{enumerate}
\end{theorem}

\textbf{Why newforms matter}: Newforms are Hecke eigenforms with well-defined arithmetic meaning. Oldforms are linear combinations of newforms from lower levels---they don't give new arithmetic information, just redundancy. To study the arithmetic of modular forms, we focus on newforms at their minimal level (conductor).

\textbf{Example}: Consider the newform $f$ associated to the elliptic curve 11a1. This newform has conductor 11, meaning:
\begin{itemize}
\item $f \in S_2^{\text{new}}(\Gamma_0(11))$ (newform at level 11)
\item $f \in S_2(\Gamma_0(22))$ (oldform at level 22, since $11 | 22$)
\item $f \in S_2(\Gamma_0(33))$ (oldform at level 33, since $11 | 33$)
\item $f \in S_2(\Gamma_0(11k))$ for any $k \geq 1$ (oldform when $k > 1$)
\end{itemize}
But we say $f$ ``has level 11'' or ``has conductor 11'' because that's where it's new.

\begin{theorem}[Modularity Theorem, Wiles-Taylor-Breuil-Conrad-Diamond]
Every elliptic curve $E/\mathbb{Q}$ of conductor $N$ corresponds to a weight 2 newform $f_E \in S_2^{\text{new}}(\Gamma_0(N))$.
\end{theorem}

\textbf{Unpacking the statement}: 
\begin{itemize}
\item The conductor of $E$ (defined geometrically via bad reduction: $N = \prod_{p \text{ bad}} p^{f_p}$)
\item \textbf{Equals} the conductor of $f_E$ (defined as the minimal level where $f_E$ appears as a newform)
\item This equality is the content of the Modularity Theorem---two a priori different notions of ``conductor'' coincide
\end{itemize}

\textbf{The correspondence}: For primes $p \nmid N$, the $p$-th Fourier coefficient of $f_E = \sum_{n=1}^{\infty} a_n q^n$ (where $q = e^{2\pi iz}$) satisfies:
$$a_p = p + 1 - \#E(\mathbb{F}_p)$$
This connects the analytic object $f_E$ to the arithmetic of $E$.

For elliptic curves, the weight is always $k = 2$. For Ramanujan's $\Delta$ function:
$$\Delta(z) = q \prod_{n=1}^{\infty} (1-q^n)^{24} = \sum_{n=1}^{\infty} \tau(n) q^n$$
we have $k = 12$. The function $\Delta$ generates the 1-dimensional space $S_{12}(\text{SL}_2(\mathbb{Z}))$.

\textbf{Why weight matters}: The weight determines transformation properties and growth conditions. Higher weight means faster decay at cusps. Weight 2 is special: it corresponds to differential forms on modular curves, linking to algebraic geometry.

\subsection{Column C: Torsion Order}

\begin{definition}
A point $P \in E(\mathbb{Q})$ is a \textbf{torsion point} if it has finite order, i.e., there exists $n \geq 1$ such that $[n]P = O$.
\end{definition}

The set of all torsion points forms a subgroup $E(\mathbb{Q})_{\text{tors}} \subseteq E(\mathbb{Q})$.

The Mordell-Weil theorem states that $E(\mathbb{Q})$ is a finitely generated abelian group. By the structure theorem:
$$E(\mathbb{Q}) \cong \mathbb{Z}^r \oplus E(\mathbb{Q})_{\text{tors}}$$
where $r \geq 0$ is the \textbf{rank} (Column E) and $E(\mathbb{Q})_{\text{tors}}$ is the finite \textbf{torsion subgroup}.

Column C shows $|E(\mathbb{Q})_{\text{tors}}|$, the order of this torsion subgroup.

\begin{theorem}[Mazur, 1977]
For elliptic curves over $\mathbb{Q}$, the torsion subgroup is one of exactly 15 possibilities:
$$\mathbb{Z}/m\mathbb{Z}, \quad m \in \{1,2,3,4,5,6,7,8,9,10,12\}$$
or
$$\mathbb{Z}/2\mathbb{Z} \times \mathbb{Z}/2m\mathbb{Z}, \quad m \in \{1,2,3,4\}$$
\end{theorem}

This is a remarkable rigidity theorem---over other number fields, many more structures are possible.

\subsection{Column D: Torsion Structure}

This gives the explicit group structure as a product of cyclic groups $\mathbb{Z}/n_1\mathbb{Z} \times \cdots \times \mathbb{Z}/n_k\mathbb{Z}$ where $n_i | n_{i+1}$ (invariant factors form).

Examples from the table:
\begin{itemize}
\item ``$\mathbb{Z}/5\mathbb{Z}$'': Cyclic group of order 5
\item ``$\mathbb{Z}/2\mathbb{Z} \times \mathbb{Z}/4\mathbb{Z}$'': Non-cyclic group of order 8
\item ``0'' or ``trivial'': Trivial group (no torsion)
\end{itemize}

\textbf{Why torsion matters}: Torsion points are defined over $\mathbb{Q}$, so they're ``easy'' rational points. For cryptography, we want large rank (many generators) but small torsion. For Diophantine problems, torsion can obstruct or facilitate solutions.

\subsection{Column E: Rank}

The \textbf{rank} $r$ is the $\mathbb{Z}$-rank of the free part of $E(\mathbb{Q})$. Geometrically, it's the number of independent rational points of infinite order.

\begin{example}
The curve $\texttt{37a1}: y^2 + y = x^3 - x$ has rank 1. One generator is $P = (0, 0)$, and every rational point is of the form $[n]P + T$ where $T$ is a torsion point and $n \in \mathbb{Z}$.
\end{example}

\textbf{The rank problem}: Computing the rank is \textbf{hard}. The best algorithms (descent methods) work provably for small rank but can fail for larger rank. The Birch and Swinnerton-Dyer (BSD) conjecture predicts:
$$\text{rank} = \text{ord}_{s=1} L(E, s)$$
where $L(E, s)$ is the L-function (see Section 6.6). For rank 0 and 1, BSD is proven in many cases, making rank computation feasible.

In our dataset ($N \leq 54$), we have:
\begin{itemize}
\item Rank 0: 138 curves
\item Rank 1: 3 curves (37a1, 43a1, 53a1)
\item Rank $\geq 2$: 0 curves
\end{itemize}

\section{Table 2: Mordell-Weil Generators and Galois Representations}

\subsection{Column B: Number of Generators}

For a rank $r$ curve, we need exactly $r$ points to generate the free part $\mathbb{Z}^r$ of $E(\mathbb{Q})/E(\mathbb{Q})_{\text{tors}}$. Column B shows this count. For rank 0 curves, this is 0. For the three rank 1 curves, this is 1.

\subsection{Column C: Generator Heights}

Given a point $P = (x, y) \in E(\mathbb{Q})$ with $x = n/d$ in lowest terms, the \textbf{naive height} is:
$$h(P) = \log \max(|n|, |d|)$$

This measures arithmetic complexity but isn't well-behaved under the group law. The \textbf{canonical height} (or Néron-Tate height) is:
$$\hat{h}(P) = \lim_{n \to \infty} \frac{h([2^n]P)}{4^n}$$

The canonical height is a \textbf{positive definite quadratic form}:
\begin{itemize}
\item $\hat{h}(P) \geq 0$ with equality iff $P \in E(\mathbb{Q})_{\text{tors}}$
\item $\hat{h}([m]P) = m^2 \hat{h}(P)$
\item $\hat{h}(P+Q) + \hat{h}(P-Q) = 2\hat{h}(P) + 2\hat{h}(Q)$ (parallelogram law)
\end{itemize}

\textbf{Quadratic form interpretation}: The parallelogram law makes $\hat{h}$ a quadratic form. The associated symmetric bilinear pairing is:
$$\langle P, Q \rangle := \frac{1}{2}(\hat{h}(P+Q) - \hat{h}(P) - \hat{h}(Q))$$

For generators $P_1, \ldots, P_r$ of $E(\mathbb{Q})/E(\mathbb{Q})_{\text{tors}}$, write any point as $P = \sum_{i=1}^r n_i P_i$ with $n_i \in \mathbb{Z}$. Then:
$$\hat{h}(P) = \sum_{i,j=1}^r n_i n_j \langle P_i, P_j \rangle$$

This is a homogeneous quadratic polynomial in $(n_1, \ldots, n_r)$ with coefficients $\langle P_i, P_j \rangle$. The matrix $H_{ij} = \langle P_i, P_j \rangle$ is positive definite.

Column C lists $\hat{h}(P_i)$ for each generator $P_i$. These are numerical approximations (typically accurate to $10^{-10}$ or better) computed using a finite approximation to the limiting process.

\textbf{Why heights matter}: Heights measure arithmetic complexity. Small height means ``simple'' rational point. Height theory connects to:
\begin{itemize}
\item Counting rational points of bounded height
\item Effective versions of Mordell's theorem
\item The regulator (next column)
\end{itemize}

\subsection{Column D: Regulator}

For generators $P_1, \ldots, P_r$ of $E(\mathbb{Q})/E(\mathbb{Q})_{\text{tors}}$, the \textbf{regulator} is:
$$\text{Reg}(E/\mathbb{Q}) = \det\left(\langle P_i, P_j \rangle\right)_{1 \leq i,j \leq r}$$

For rank 1, the regulator equals the height of the generator: $\text{Reg} = \hat{h}(P_1)$. For rank 0, we set $\text{Reg} = 1$ by convention.

\textbf{Why the regulator matters}: It appears in the BSD conjecture:
$$\lim_{s \to 1} \frac{L(E, s)}{(s-1)^r} = \frac{\Omega_E \cdot \text{Reg}(E/\mathbb{Q}) \cdot \prod_p c_p \cdot |\text{Sha}(E/\mathbb{Q})|}{|E(\mathbb{Q})_{\text{tors}}|^2}$$
where $\Omega_E$ is a period, $c_p$ are Tamagawa numbers, and $\text{Sha}$ is the Tate-Shafarevich group. The regulator measures the ``volume'' of the lattice of rational points.

\subsection{Column E: Galois Representation Surjective for All Primes?}

Now we enter representation theory. For each prime $\ell$, consider the $\ell$-torsion:
$$E[\ell] = \{P \in E(\overline{\mathbb{Q}}) : [\ell]P = O\}$$

This is a 2-dimensional vector space over $\mathbb{F}_\ell$: $E[\ell] \cong \mathbb{F}_\ell^2$ as abelian groups.

The absolute Galois group $G_{\mathbb{Q}} = \text{Gal}(\overline{\mathbb{Q}}/\mathbb{Q})$ acts on $E[\ell]$ by:
$$\sigma \cdot P = \sigma(x, y) \text{ for } P = (x, y) \in E[\ell]$$
where $\sigma$ extends to coordinates. This action respects the group structure, giving a representation:
$$\rho_{E,\ell}: G_{\mathbb{Q}} \to \text{Aut}(E[\ell]) \cong \text{GL}_2(\mathbb{F}_\ell)$$

\textbf{The question}: Is $\rho_{E,\ell}$ surjective onto $\text{GL}_2(\mathbb{F}_\ell)$?

Column E answers: ``Is $\rho_{E,\ell}$ surjective for \emph{all} primes $\ell$?'' 
\begin{itemize}
\item Y (Yes/True) = surjective for all $\ell$
\item N (No/False) = non-surjective for at least one $\ell$
\end{itemize}

\textbf{What does Galois action mean?} The Galois group $G_{\mathbb{Q}}$ permutes the algebraic numbers. When it acts on $E[\ell]$, it permutes the $\ell$-torsion points while preserving the group structure. Larger image of $\rho_{E,\ell}$ means Galois ``mixes'' the $\ell$-torsion points more thoroughly.

\textbf{How is ``more'' quantified?} By the size $|\text{im}(\rho_{E,\ell})|$ and structure of the image subgroup. Maximal mixing occurs when $\text{im}(\rho_{E,\ell}) = \text{GL}_2(\mathbb{F}_\ell)$, which has order:
$$|\text{GL}_2(\mathbb{F}_\ell)| = \ell(\ell-1)^2(\ell+1)$$

\textbf{Why this matters}: Galois representations encode arithmetic information:
\begin{itemize}
\item \textbf{Local data}: For $\ell \nmid N$, the trace $\text{tr}(\rho_{E,\ell}(\text{Frob}_p))$ equals $a_p$, the coefficient in $f_E = \sum a_n q^n$
\item \textbf{Torsion fields}: $\mathbb{Q}(E[\ell])$ is the fixed field of $\ker(\rho_{E,\ell})$
\item \textbf{Modularity}: The representation connects to automorphic forms
\item \textbf{L-functions}: Functional equations and special values
\end{itemize}

\begin{definition}
An elliptic curve $E/\mathbb{Q}$ has \textbf{complex multiplication} (CM) if its endomorphism ring over $\overline{\mathbb{Q}}$ is strictly larger than $\mathbb{Z}$:
$$\text{End}_{\overline{\mathbb{Q}}}(E) \supsetneq \mathbb{Z}$$
By theory, when this occurs, $\text{End}_{\overline{\mathbb{Q}}}(E) \cong \mathcal{O}_K$ is an order in an imaginary quadratic field $K = \mathbb{Q}(\sqrt{D})$ with $D < 0$.
\end{definition}

\begin{theorem}[Serre]
For non-CM curves over $\mathbb{Q}$, the representation $\rho_{E,\ell}$ is surjective for all but finitely many primes $\ell$.
\end{theorem}

Most entries in Column E are ``N'' because even one exceptional prime makes it false. Table 3 lists these exceptional primes.

\section{Table 3: Galois Exceptional Primes and Modular Data}

\subsection{Column B: Galois Representation Exceptional Primes}

This lists the finite set of primes $\ell$ where $\rho_{E,\ell}$ fails to be surjective. Notation: ``2|3|5'' means primes 2, 3, and 5.

\textbf{Why non-surjectivity occurs}:

To understand this, we need the subgroup structure of $\text{GL}_2(\mathbb{F}_\ell)$.

\begin{definition}
The \textbf{Borel subgroup} $B \subseteq \text{GL}_2(\mathbb{F}_\ell)$ is the subgroup of upper-triangular matrices:
$$B = \left\{ \begin{pmatrix} a & b \\ 0 & d \end{pmatrix} : ad \neq 0 \right\}$$
\end{definition}

\textbf{Interpretation}: $B$ is the stabilizer of the line $\langle e_1 \rangle \subseteq \mathbb{F}_\ell^2$ (spanned by the first basis vector). Elements of $B$ scale this line and act arbitrarily on the quotient space.

\begin{definition}
A \textbf{Cartan subgroup} of $\text{GL}_2(\mathbb{F}_\ell)$ is a maximal abelian subgroup. There are two conjugacy classes:

\textbf{Split Cartan}: The diagonal matrices
$$C_s = \left\{ \begin{pmatrix} a & 0 \\ 0 & d \end{pmatrix} : ad \neq 0 \right\} \cong \mathbb{F}_\ell^* \times \mathbb{F}_\ell^*$$
This has order $(\ell-1)^2$.

\textbf{Non-split Cartan}: The image of the embedding $\mathbb{F}_{\ell^2}^* \hookrightarrow \text{GL}_2(\mathbb{F}_\ell)$, where multiplication by elements of $\mathbb{F}_{\ell^2}$ acts as $\mathbb{F}_\ell$-linear maps on $\mathbb{F}_{\ell^2} \cong \mathbb{F}_\ell^2$. This has order $\ell^2 - 1$ and is cyclic.
\end{definition}

\textbf{Explicit non-split Cartan}: If $\mathbb{F}_{\ell^2} = \mathbb{F}_\ell[t]/(t^2 - d)$ where $d \in \mathbb{F}_\ell$ is not a square, then the element $t$ acts as:
$$t \mapsto \begin{pmatrix} 0 & d \\ 1 & 0 \end{pmatrix}$$
and generates (with scalars) a cyclic subgroup of order $\ell^2-1$.

Now we can describe when $\rho_{E,\ell}$ is not surjective:

\begin{enumerate}
\item \textbf{Rational $\ell$-torsion}: If $E(\mathbb{Q})[\ell] \neq 0$, then $\rho_{E,\ell}$ fixes a non-zero point $P \in E[\ell]$, so the image is contained in the stabilizer of $\langle P \rangle$, which is conjugate to a Borel subgroup.

Example: Curves 11a1, 11a2, 11a3 all have 5-torsion over $\mathbb{Q}$ (visible in Column D of Table 1), so $\rho_{E,5}$ is not surjective.

\item \textbf{Rational $\ell$-isogeny}: If there's an isogeny $E \to E'$ of degree $\ell$ defined over $\mathbb{Q}$, the image of $\rho_{E,\ell}$ is contained in the normalizer of a Cartan subgroup.

\item \textbf{Complex multiplication}: If $E$ has CM, the extra endomorphisms force the image of $\rho_{E,\ell}$ to lie in the normalizer of a Cartan subgroup for all $\ell$.
\end{enumerate}

For small $\ell$ (especially $\ell = 2, 3$), there can be \textbf{exceptional subgroups}---subgroups that don't fit the Borel/Cartan pattern. For instance, $\text{GL}_2(\mathbb{F}_2) \cong S_3$ (order 6) has a richer subgroup lattice than larger cases. For $\ell \geq 5$ and non-CM curves, Serre's analysis shows the image is typically all of $\text{GL}_2(\mathbb{F}_\ell)$, with exceptions coming from the cases above.

\subsection{Column C: Galois Image Type}

A classification of the Galois representation:
\begin{itemize}
\item ``surjective'': $\text{im}(\rho_{E,\ell}) = \text{GL}_2(\mathbb{F}_\ell)$ for all $\ell$
\item ``non-surjective at [list]'': Lists exceptional primes
\item ``CM'': Curve has complex multiplication
\end{itemize}

CM curves are special:
\begin{itemize}
\item Over $\mathbb{Q}$, there are only 13 possible CM $j$-invariants, corresponding to imaginary quadratic fields of class number 1: discriminants $D \in \{-3, -4, -7, -8, -11, -12, -16, -19, -27, -28, -43, -67, -163\}$
\item Galois representations are never surjective (image lies in normalizer of Cartan)
\item L-functions have special values (related to Hecke characters)
\item Heights are easier to compute
\end{itemize}

Examples in our table: Curves 27a1-27a4 have CM by $\mathbb{Z}[\omega]$ where $\omega = e^{2\pi i/3}$ (discriminant $-3$).

\subsection{Column D: Index of $\Gamma_0(N)$ in $\text{SL}_2(\mathbb{Z})$}

The \textbf{congruence subgroup} $\Gamma_0(N) \subset \text{SL}_2(\mathbb{Z})$ consists of matrices:
$$\Gamma_0(N) = \left\{ \begin{pmatrix} a & b \\ c & d \end{pmatrix} \in \text{SL}_2(\mathbb{Z}) : c \equiv 0 \pmod{N} \right\}$$

The index is:
$$[\text{SL}_2(\mathbb{Z}) : \Gamma_0(N)] = N \prod_{p | N} \left(1 + \frac{1}{p}\right)$$

Examples:
\begin{itemize}
\item $N = 11$ (prime): Index = $11 \cdot \frac{12}{11} = 12$
\item $N = 15 = 3 \cdot 5$: Index = $15 \cdot \frac{4}{3} \cdot \frac{6}{5} = 24$
\item $N = 1$ ($\Delta$): Index = 1 (since $\Gamma_0(1) = \text{SL}_2(\mathbb{Z})$)
\end{itemize}

\textbf{Why this matters}:

\textbf{1. Modular forms perspective}: The dimension of the space of weight $k$ cusp forms grows like:
$$\dim S_k(\Gamma_0(N)) \sim \frac{k-1}{12} [\text{SL}_2(\mathbb{Z}) : \Gamma_0(N)]$$
for large $N$. Larger index means larger space of modular forms.

\textbf{2. Representation theory}: The group $\Gamma_0(N)$ acts on the upper half-plane $\mathbb{H}$, and modular forms are functions on $\Gamma_0(N) \backslash \mathbb{H}$ with nice transformation properties. The index tells us:
\begin{itemize}
\item Size of the fundamental domain
\item Number of coset representatives
\item Complexity of Hecke operators
\end{itemize}

\textbf{3. Modularity}: The curve $E$ corresponds to a newform $f \in S_2(\Gamma_0(N))$. This newform is an eigenfunction for Hecke operators, and its Fourier coefficients $a_n$ encode arithmetic data about $E$. The level $N$ and index control:
\begin{itemize}
\item Which Hecke algebras act
\item Congruences between modular forms
\item Galois representations attached to $f$
\end{itemize}

\section{Table 4: Modular Curves and L-functions}

\subsection{Column B: Genus of $X_0(N)$}

The modular curve $X_0(N)$ is defined as:
$$X_0(N) = \Gamma_0(N) \backslash (\mathbb{H} \cup \mathbb{P}^1(\mathbb{Q}))$$
the compactification of the quotient of the upper half-plane by $\Gamma_0(N)$.

As a Riemann surface, $X_0(N)$ has genus $g$ computed via Riemann-Hurwitz:
$$g = 1 + \frac{[\text{SL}_2(\mathbb{Z}):\Gamma_0(N)]}{12} - \frac{\nu_2}{4} - \frac{\nu_3}{3} - \frac{\nu_\infty}{2}$$
where $\nu_2, \nu_3$ count elliptic points and $\nu_\infty$ counts cusps.

\textbf{Why genus matters}:

\textbf{1. Dimension of modular forms}: For weight 2:
$$\dim S_2(\Gamma_0(N)) = g \text{ for } N > 1$$

Example: $X_0(37)$ has genus 2, so there are 2 independent weight-2 newforms. These correspond to curves 37a1 and 37b1.

\textbf{2. Rational points}: The curve $X_0(N)$ itself is an algebraic curve over $\mathbb{Q}$. Its rational points $X_0(N)(\mathbb{Q})$ parametrize pairs $(E, C)$ where:
\begin{itemize}
\item $E$ is an elliptic curve over $\mathbb{Q}$
\item $C \subset E$ is a cyclic subgroup of order $N$
\end{itemize}

By Mazur's theorem, for $N > 163$, we have $X_0(N)(\mathbb{Q}) = \{\text{cusps}\}$, so there are no non-cuspidal rational points.

\textbf{3. Genus and complexity}:
\begin{itemize}
\item Genus 0: Rational curve (e.g., $X_0(1), X_0(2), \ldots, X_0(10)$)
\item Genus 1: Elliptic curve (e.g., $X_0(11), X_0(14), X_0(15)$)
\item Genus $\geq 2$: General type, Mordell conjecture (Faltings) applies
\end{itemize}

\subsection{Column C: Number of Cusps}

Cusps are points added in compactification, corresponding to $\Gamma_0(N)$-orbits in $\mathbb{P}^1(\mathbb{Q})$. The count is:
$$\#\{\text{cusps}\} = \sum_{d | N} \phi(\gcd(d, N/d))$$
where $\phi$ is Euler's totient function.

Examples:
\begin{itemize}
\item $N = 11$ (prime): 2 cusps
\item $N = 30 = 2 \cdot 3 \cdot 5$: 12 cusps
\item $N = 1$: 1 cusp (the point at $\infty$)
\end{itemize}

\textbf{Why cusps matter}:

\textbf{1. Fourier expansions}: Modular forms have Fourier expansions at each cusp. A cusp form $f \in S_k(\Gamma_0(N))$ vanishes at all cusps, which imposes strong constraints.

\textbf{2. Rational points}: Cusps are always rational points of $X_0(N)$, representing ``degenerate'' elliptic curves (nodal or cuspidal cubics).

\textbf{3. Dimension formulas}: The genus formula involves the number of cusps, connecting topology to arithmetic.

\subsection{Column D: Has Complex Multiplication?}

As defined in Section 4.5, a curve has CM if $\text{End}_{\overline{\mathbb{Q}}}(E) \supsetneq \mathbb{Z}$.

\textbf{Why modular form workers care about CM}:

\textbf{1. Special L-values}: For CM curves, the L-function $L(E, s)$ factors as:
$$L(E, s) = L(\psi, s)$$
where $\psi$ is a Hecke character of the imaginary quadratic field. This connects to:
\begin{itemize}
\item Class field theory
\item Algebraic Hecke characters
\item Explicit formulas for special values
\end{itemize}

\textbf{2. Congruences}: CM forms often satisfy congruences with non-CM forms modulo small primes, leading to:
\begin{itemize}
\item Connections between different L-functions
\item Iwasawa theory
\item $p$-adic L-functions
\end{itemize}

\textbf{3. Explicit construction}: CM curves can be constructed explicitly via the theory of complex multiplication:
\begin{itemize}
\item Given discriminant $D$, construct $j$-invariant using modular functions
\item Reduces to class field theory of imaginary quadratic fields
\item Enables explicit computation of Fourier coefficients
\end{itemize}

\subsection{Column E: CM Discriminant}

If $E$ has CM, this is the discriminant $D < 0$ of the imaginary quadratic order $\text{End}(E)$. Standard values: $D \in \{-3, -4, -7, -8, -11, -12, -16, \ldots\}$.

Examples from table:
\begin{itemize}
\item Curves 27a1-27a4: $D = -3$ (CM by $\mathbb{Z}[\omega]$, $\omega^3 = 1$)
\item Curves 32a1-32a4: $D = -4$ (CM by $\mathbb{Z}[i]$, $i^2 = -1$)
\item Curves 49a1-49a4: $D = -7$ (CM by $\mathbb{Z}[\frac{1+\sqrt{-7}}{2}]$)
\end{itemize}

\textbf{Connection to class field theory}: The field $K = \mathbb{Q}(\sqrt{D})$ has class number $h_D$. The 13 discriminants mentioned earlier ($D = -3, -4, -7, -8, -11, -12, -16, -19, -27, -28, -43, -67, -163$) are exactly those with $h_D = 1$ (class number one). For these discriminants, the CM elliptic curves have rational $j$-invariants.

\subsection{Column F: Root Number}

The L-function of $E$ is:
$$L(E, s) = \prod_{p \nmid N} \frac{1}{1 - a_p p^{-s} + p^{1-2s}} \cdot \prod_{p | N} L_p(E, p^{-s})$$
where the second product accounts for bad primes.

This has an analytic continuation to all $s \in \mathbb{C}$ and satisfies a functional equation:
$$\Lambda(E, s) = w \cdot \Lambda(E, 2-s)$$
where $\Lambda(E, s) = N^{s/2}(2\pi)^{-s}\Gamma(s) L(E, s)$ is the completed L-function, and $w \in \{+1, -1\}$ is the \textbf{root number}.

\textbf{Why the root number matters}:

\textbf{1. Birch and Swinnerton-Dyer conjecture}: The BSD conjecture predicts:
$$w = (-1)^{r}$$
where $r$ is the rank. Thus:
\begin{itemize}
\item $w = +1$ predicts even rank (including rank 0)
\item $w = -1$ predicts odd rank (including rank 1)
\end{itemize}

In our table:
\begin{itemize}
\item 138 curves have $w = +1$ and rank 0 (consistent with BSD)
\item 3 curves (37a1, 43a1, 53a1) have $w = -1$ and rank 1 (consistent with BSD)
\end{itemize}

\textbf{2. Representation theory}: The root number is a product of local signs:
$$w = \prod_{p \leq \infty} w_p$$
where each local factor $w_p = \pm 1$ encodes arithmetic data at $p$. For finite primes $p | N$:
\begin{itemize}
\item \textbf{Multiplicative reduction at $p$}: Typically $w_p = -1$ (``bad for parity'')
\item \textbf{Additive reduction at $p$}: Typically $w_p = +1$ (``good for parity'')
\end{itemize}

\textbf{What does ``good/bad for parity'' mean?} Since $w = \prod_p w_p$ and we want to know the sign:
\begin{itemize}
\item Each $w_p = -1$ contributes a negative sign to the product
\item If there's an \emph{odd} number of primes with $w_p = -1$, then $w = -1$ (predicting odd rank)
\item If there's an \emph{even} number of primes with $w_p = -1$, then $w = +1$ (predicting even rank)
\end{itemize}

So ``bad for parity'' means $w_p = -1$ flips the sign, potentially changing even rank prediction to odd (or vice versa). ``Good for parity'' means $w_p = +1$ doesn't flip the sign.

\textbf{3. Modular forms}: The root number relates to the Atkin-Lehner involution $W_N$ acting on $S_2(\Gamma_0(N))$:
$$W_N: f(z) \mapsto \frac{1}{\sqrt{N}} f(-1/(Nz))$$

The newform $f_E$ is an eigenfunction: $W_N(f_E) = w \cdot f_E$.

\textbf{4. Deep connections}: Via the Modularity Theorem:
\begin{itemize}
\item Root number connects to Galois representations (via local Galois data)
\item Relates to epsilon factors in Langlands program
\item Appears in formulations of equivariant Tamagawa number conjectures
\end{itemize}

\section{Why Representation Theory Matters: A Summary}

Throughout these tables, representation-theoretic data appears repeatedly. Why do modular form workers care so much about Galois representations?

\subsection{The Langlands Program}

The overarching framework is the \textbf{Langlands program}, which proposes deep connections between:
\begin{itemize}
\item \textbf{Automorphic side}: Modular forms, automorphic representations
\item \textbf{Galois side}: Galois representations, motives
\end{itemize}

For elliptic curves, the Modularity Theorem is a special case: the curve $E$ (geometric object with Galois action) corresponds to the newform $f_E$ (analytic object with Hecke action).

The Galois representation $\rho_{E,\ell}$ encodes how Galois acts on $\ell$-torsion, while the modular form $f_E$ encodes the same information through Fourier coefficients. The correspondence is:
$$\text{tr}(\rho_{E,\ell}(\text{Frob}_p)) \equiv a_p \pmod{\ell}$$
for primes $p \nmid \ell N$.

\subsection{Local-Global Principle}

Galois representations encode:
\begin{itemize}
\item \textbf{Global data}: Via $\rho_{E,\ell}: G_{\mathbb{Q}} \to \text{GL}_2(\mathbb{F}_\ell)$
\item \textbf{Local data}: Via restrictions $\rho_{E,\ell}|_{G_{\mathbb{Q}_p}}$ for each prime $p$
\end{itemize}

The local data determines:
\begin{itemize}
\item Reduction type at $p$ (good, multiplicative, additive)
\item Conductor exponent $f_p$ at $p$
\item Local root number $w_p$
\end{itemize}

Modular forms see this same data through Fourier coefficients and Hecke eigenvalues. The representation-theoretic formulation unifies these perspectives.

\subsection{Congruences and $p$-adic Theory}

When two modular forms have congruent Fourier coefficients modulo a prime $p$, their associated Galois representations are congruent modulo $p$. This leads to:
\begin{itemize}
\item \textbf{Iwasawa theory}: Growth of Selmer groups in $\mathbb{Z}_p$-extensions
\item \textbf{$p$-adic L-functions}: Interpolating special values $p$-adically
\item \textbf{Bloch-Kato conjectures}: Relating regulators to $L$-values via $K$-theory
\end{itemize}

Understanding the Galois representation modulo $p$ is essential for these deep arithmetic applications.

\subsection{Modularity Lifting}

The proof of Fermat's Last Theorem (Wiles-Taylor) used \textbf{modularity lifting theorems}: given a Galois representation $\rho$ that looks modular mod $p$ (matches a modular form mod $p$), prove $\rho$ is modular (comes from a modular form over characteristic 0). This requires:
\begin{itemize}
\item Detailed control over Galois cohomology $H^1(G_{\mathbb{Q}}, \text{ad}^0(\rho))$
\item Deformation theory of representations
\item Understanding ramification and local conditions
\end{itemize}

The representation-theoretic data in Tables 2 and 3---exceptional primes, image structure, CM properties---is exactly the input needed for such theorems. To prove modularity of a Galois representation, one must understand:
\begin{itemize}
\item At which primes is it ramified? (Conductor)
\item What is the local structure at ramified primes? (Reduction type)
\item What are the irreducibility properties? (Image of $\rho$)
\item Are there congruences with known modular forms? (Exceptional primes)
\end{itemize}

\section{Conclusion}

These four tables encode a rich tapestry of connections between:
\begin{itemize}
\item Algebraic geometry (elliptic curves as varieties)
\item Group theory (rational points as abelian groups)
\item Complex analysis (modular forms as holomorphic functions)
\item Representation theory (Galois representations)
\item Number theory (L-functions, BSD conjecture)
\end{itemize}

For researchers in modular forms, the representation-theoretic perspective is not optional---it's the language in which modern number theory is written. The Galois representations provide:
\begin{itemize}
\item A unified framework (Langlands program)
\item Computational tools (modularity, congruences)
\item Deep conjectures (BSD, Bloch-Kato)
\item Connections to other areas (arithmetic geometry, algebraic topology)
\end{itemize}

Even for those primarily interested in modular forms as analytic objects, understanding the associated Galois representations is essential for grasping the deeper arithmetic meaning of Fourier coefficients, Hecke eigenvalues, and L-functions.

The data in these tables---conductor exponents, torsion structures, generator heights, exceptional primes, genus, root numbers---are not isolated curiosities but fundamental invariants that:
\begin{itemize}
\item Determine the structure of Galois representations
\item Control dimensions of modular form spaces
\item Predict ranks via the BSD conjecture
\item Enable explicit computations of L-values
\item Guide the search for rational points
\end{itemize}

Understanding what each column means, and why it matters, is the first step toward appreciating the profound unity underlying the arithmetic of elliptic curves.

\end{document}
